%==============================================================================
%  Eletromagnetismo — Uma breve introdução
%  Notas transcritas e reorganizadas a partir de manuscrito.
%  Base bibliográfica: A. Zangwill, "Modern Electrodynamics".
%
%  >>> COMO ALTERAR A FORMATAÇÃO <<<
%  Todos os parâmetros de formatação (fonte, margens, espaçamento, cor das
%  marcações) estão reunidos e comentados no bloco "PARÂMETROS" logo abaixo.
%  Basta editar ali; o corpo do texto não precisa ser tocado.
%
%  Convenção de revisão:
%  - Tudo em PRETO é transcrição fiel do manuscrito.
%  - Tudo em VERMELHO foi ADICIONADO ou CORRIGIDO por quem digitou.
%    Trechos curtos:  \add{...}
%    Trechos longos:  ambiente  \begin{nota} ... \end{nota}
%    (ambos começam com o adereço vermelho  >  e terminam em vermelho)
%==============================================================================
\documentclass[12pt,a4paper]{article}

%------------------------------------------------------------------ codificação
\usepackage[T1]{fontenc}          % fontes Computer Modern com acentuação (EC/T1)
\usepackage[utf8]{inputenc}
\usepackage{babel}
\babelprovide[import, main]{brazilian}   % português do Brasil (hifenização/rótulos)

%------------------------------------------------------------------ matemática
\usepackage{amsmath,amssymb,mathtools}
\usepackage{bm}                   % negrito matemático para vetores

%==============================================================================
%  PARÂMETROS  (edite aqui para reformatar tudo)
%==============================================================================
% -- Fonte -------------------------------------------------------------------
%    Padrão: Computer Modern (fonte nativa do LaTeX), conforme solicitado.
%    Para trocar por Latin Modern, descomente a linha abaixo (se instalado):
% \usepackage{lmodern}
%    Para Times/ABNT estrito: \usepackage{mathptmx}

% -- Margens (ABNT NBR 14724: esq./sup. 3 cm, dir./inf. 2 cm) -----------------
\usepackage[a4paper,left=3cm,top=3cm,right=2cm,bottom=2cm]{geometry}

% -- Espaçamento entre linhas (ABNT: 1,5) ------------------------------------
\usepackage{setspace}
\onehalfspacing

% -- Recuo de primeira linha (ABNT: parágrafo indentado, inclusive o 1.o) -----
\usepackage{indentfirst}
\setlength{\parindent}{1.25cm}

% -- Cores e marcação de revisão ---------------------------------------------
\usepackage{xcolor}
\definecolor{coradd}{RGB}{200,0,0}       % vermelho das adições/correções
\newcommand{\adereco}{\textcolor{coradd}{\ensuremath{\blacktriangleright}}} % adereço
% adição/correção curta (em linha): abre com adereço, texto vermelho
\newcommand{\add}[1]{\adereco\,\textcolor{coradd}{#1}}
% adição/correção longa (bloco): começa com adereço e termina em vermelho
\newenvironment{nota}
  {\par\nobreak\vspace{2pt}\noindent\adereco\ \bgroup\color{coradd}}
  {\ \textcolor{coradd}{$\blacktriangleleft$}\egroup\par\vspace{2pt}}

% -- Tolerância de justificação (evita texto invadindo a margem) --------------
\setlength{\emergencystretch}{3em}

% -- Links e utilidades -------------------------------------------------------
\usepackage{amsthm}
\usepackage[hidelinks]{hyperref}

% -- Ambientes de destaque ----------------------------------------------------
\theoremstyle{definition}
\newtheorem{exemplo}{Exemplo}[section]
\theoremstyle{remark}
\newtheorem*{obs}{Observação}

%------------------------------------------------------------------ atalhos math
\renewcommand{\vec}[1]{\bm{#1}}          % vetores em negrito
\newcommand{\uvec}[1]{\hat{\bm{#1}}}     % versores
\newcommand{\grad}{\vec{\nabla}}
\newcommand{\dpar}[2]{\dfrac{\partial #1}{\partial #2}}
\newcommand{\rsep}{\bm{\mathcal{R}}}     % vetor separação (script-r de Zangwill)
\newcommand{\usep}{\hat{\bm{\mathcal{R}}}}
\newcommand{\Rsep}{\mathcal{R}}          % módulo do vetor separação
\newcommand{\eps}{\varepsilon_{0}}

%==============================================================================
\begin{document}

%------------------------------------------------------------------ título ABNT
\begin{center}
  {\Large\bfseries Eletromagnetismo — Uma breve introdução}\\[4pt]
  {\large Preliminares Matemáticos e Eletrostática}\\[10pt]
  {\normalsize Notas de estudo}\\[2pt]
  \add{Baseado em A.~Zangwill, \textit{Modern Electrodynamics} (unidades SI).}\\[6pt]
  \rule{\textwidth}{0.4pt}
\end{center}

\begin{nota}
Legenda desta revisão: o texto em preto reproduz o manuscrito; o texto em
vermelho, iniciado pelo adereço $\blacktriangleright$, indica passagens
adicionadas ou corrigidas na digitação (fechamento das lacunas de notação,
resultados finais de contas apenas esboçadas e acertos de sinais/índices).
\end{nota}

\tableofcontents
\bigskip
\hrule
\bigskip

%##############################################################################
\section{Preliminares matemáticos}
%##############################################################################

\subsection{Álgebra vetorial e a convenção de Einstein}

Seja $\{\vec{e}_i\}$, $i=1,2,3$, um conjunto ortonormal no espaço
$(\mathbb{R}^3,\langle\,\cdot\,,\cdot\,\rangle)$, onde $\langle\,\cdot\,,\cdot\,\rangle$
denota um produto interno. Um vetor $\vec{V}$, suposto expresso nessa base, é dado
pela combinação linear
\begin{equation}
  \vec{V}=\sum_{i=1}^{3} V_i\,\vec{e}_i .
\end{equation}
Usando a \textbf{convenção de soma de Einstein} (índices repetidos, um em cima e
outro embaixo, são somados), escrevemos simplesmente $\vec{V}=V^{i}\vec{e}_i$.

Considere $\{\vec{e}\,'_j\}$ outro conjunto ortonormal. Ele se relaciona ao
primeiro por
\begin{equation}
  \vec{e}\,'_j=\sum_i A^{\,i}{}_{j}\,\vec{e}_i=A^{\,i}{}_{j}\,\vec{e}_i,
  \qquad\text{onde}\qquad
  A^{\,i}{}_{j}=\langle \vec{e}_i,\vec{e}\,'_j\rangle .
\end{equation}
\begin{nota}
A matriz $A^{\,i}{}_{j}$ é a matriz de mudança de base; como ambas as bases são
ortonormais, $A$ é ortogonal ($A^{\mathsf{T}}A=\mathbb{1}$), de modo que
preserva comprimentos e ângulos.
\end{nota}

Definimos que, em um \textbf{sistema cartesiano} $\{\vec{e}_i\}$,
\begin{equation}
  \langle \vec{e}_i,\vec{e}_j\rangle=\delta_{ij}
  \qquad\text{e}\qquad
  \dpar{\vec{e}_j}{x^i}=\vec{0}.
  \label{eq:cartesiano}
\end{equation}
Isto é, os versores cartesianos são ortonormais e constantes (não variam de ponto
a ponto). O \textbf{vetor posição} é
\begin{equation}
  \vec{r}=x\,\vec{e}_1+y\,\vec{e}_2+z\,\vec{e}_3=x^{i}\vec{e}_i .
\end{equation}
Numa segunda base, $\vec{r}=x^{i}\vec{e}_i=x^{i}A^{\,j}{}_{i}\vec{e}\,'_j
=\tilde{x}^{\,j}\vec{e}\,'_j$, com $\tilde{x}^{\,j}=A^{\,j}{}_{i}x^{i}$; as
componentes transformam-se de modo contravariante às da base.

\subsection{As quatro operações vetoriais}

Grandezas vetoriais possuem \emph{módulo} e \emph{direção}; grandezas escalares,
apenas módulo. Denotamos vetores por $\vec{A}$ e seu módulo por $|\vec{A}|=A$.
Sobre $\mathbb{R}^3$ definem-se quatro operações.

\paragraph{(i) Adição.} Comutativa e associativa,
\(
  \vec{A}+\vec{B}=\vec{B}+\vec{A}.
\)

\paragraph{(ii) Multiplicação por escalar.} Distributiva,
\(
  a(\vec{A}+\vec{B})=a\vec{A}+a\vec{B}.
\)

\paragraph{(iii) Produto escalar.}
\begin{equation}
  \vec{A}\cdot\vec{B}=AB\cos\theta,
\end{equation}
sendo $\theta$ o ângulo entre os vetores. É comutativo e distributivo:
\(
  \vec{A}\cdot\vec{B}=\vec{B}\cdot\vec{A}
\)
e
\(
  \vec{A}\cdot(\vec{B}+\vec{C})=\vec{A}\cdot\vec{B}+\vec{A}\cdot\vec{C}.
\)
Vetores perpendiculares têm $\vec{A}\cdot\vec{B}=0$; e $\vec{A}\cdot\vec{A}=A^2$.

\paragraph{(iv) Produto vetorial.}
\begin{equation}
  \vec{A}\times\vec{B}=AB\sin\theta\,\uvec{n},
\end{equation}
com $\uvec{n}$ o versor perpendicular ao plano de $\vec{A}$ e $\vec{B}$, orientado
pela regra da mão direita. É \add{anticomutativo} e distributivo:
\begin{equation}
  \vec{A}\times\vec{B}=-\,\vec{B}\times\vec{A},
  \qquad
  \vec{A}\times(\vec{B}+\vec{C})=\vec{A}\times\vec{B}+\vec{A}\times\vec{C},
  \qquad
  \vec{A}\times\vec{A}=\vec{0}.
\end{equation}

\subsection{Componentes cartesianas}

Na base ortonormal $\{\uvec{x},\uvec{y},\uvec{z}\}$, escrevemos
$\vec{A}=A_x\uvec{x}+A_y\uvec{y}+A_z\uvec{z}$. Como
$\uvec{x}\cdot\uvec{x}=\uvec{y}\cdot\uvec{y}=\uvec{z}\cdot\uvec{z}=1$ e os produtos
cruzados são nulos,
\begin{align}
  \vec{A}\cdot\vec{B}&=A_xB_x+A_yB_y+A_zB_z,\\
  \vec{A}\times\vec{B}&=
  \begin{vmatrix}
    \uvec{x} & \uvec{y} & \uvec{z}\\
    A_x & A_y & A_z\\
    B_x & B_y & B_z
  \end{vmatrix}.
\end{align}
\begin{nota}
Em notação de índices, com o símbolo de Levi-Civita $\epsilon_{ijk}$,
\(
  (\vec{A}\times\vec{B})_i=\epsilon_{ijk}A_jB_k,
\)
o que reproduz as duas identidades acima.
\end{nota}

\subsection{Produtos triplos}

\paragraph{Produto triplo escalar.} Vale a permutação cíclica
\begin{equation}
  \vec{A}\cdot(\vec{B}\times\vec{C})
  =\vec{B}\cdot(\vec{C}\times\vec{A})
  =\vec{C}\cdot(\vec{A}\times\vec{B})
  =\begin{vmatrix}
    A_x & A_y & A_z\\ B_x & B_y & B_z\\ C_x & C_y & C_z
  \end{vmatrix},
\end{equation}
\add{numericamente igual ao volume do paralelepípedo gerado por
$\vec{A},\vec{B},\vec{C}$.}

\paragraph{Produto triplo vetorial.} Regra ``BAC-CAB'':
\begin{equation}
  \vec{A}\times(\vec{B}\times\vec{C})
  =\vec{B}\,(\vec{A}\cdot\vec{C})-\vec{C}\,(\vec{A}\cdot\vec{B}).
\end{equation}

\subsection{Vetores posição, deslocamento e separação}

O vetor posição é $\vec{r}=x\uvec{x}+y\uvec{y}+z\uvec{z}$, com
$r=\sqrt{x^2+y^2+z^2}$ e $\uvec{r}=\vec{r}/r$. O deslocamento infinitesimal é
\begin{equation}
  \mathrm{d}\vec{l}=\mathrm{d}x\,\uvec{x}+\mathrm{d}y\,\uvec{y}+\mathrm{d}z\,\uvec{z}.
\end{equation}
Em problemas de fontes, distingue-se o ponto de campo $\vec{r}$ do ponto de fonte
$\vec{r}\,'$. O \textbf{vetor separação} (o \emph{script-r} de Zangwill) é
\begin{equation}
  \rsep=\vec{r}-\vec{r}\,',\qquad
  \Rsep=|\vec{r}-\vec{r}\,'|,\qquad
  \usep=\frac{\vec{r}-\vec{r}\,'}{|\vec{r}-\vec{r}\,'|}.
\end{equation}

\subsection{Cálculo diferencial: gradiente e o operador \texorpdfstring{$\grad$}{nabla}}

Para $f(x)$, a derivada mede a taxa de variação: $\mathrm{d}f=f'(x)\,\mathrm{d}x$.
Para um campo escalar $f(x,y,z)$,
\begin{equation}
  \mathrm{d}f=\dpar{f}{x}\mathrm{d}x+\dpar{f}{y}\mathrm{d}y+\dpar{f}{z}\mathrm{d}z
           =\grad f\cdot\mathrm{d}\vec{l},
\end{equation}
o que \emph{define} o \textbf{gradiente}
\begin{equation}
  \grad f=\dpar{f}{x}\uvec{x}+\dpar{f}{y}\uvec{y}+\dpar{f}{z}\uvec{z}.
\end{equation}
\begin{nota}
Interpretação: $\mathrm{d}f=|\grad f|\,|\mathrm{d}\vec{l}|\cos\theta$ é máximo
quando $\mathrm{d}\vec{l}\parallel\grad f$. Logo $\grad f$ aponta na direção de
crescimento mais rápido de $f$, e $|\grad f|$ é essa taxa máxima. Onde
$\grad f=\vec{0}$, o ponto é estacionário.
\end{nota}
O operador \textbf{del} é o objeto vetorial
\begin{equation}
  \grad=\uvec{x}\,\dpar{}{x}+\uvec{y}\,\dpar{}{y}+\uvec{z}\,\dpar{}{z},
\end{equation}
que atua de três formas: como gradiente ($\grad f$), como divergente
($\grad\cdot\vec{v}$) e como rotacional ($\grad\times\vec{v}$).

\paragraph{Divergente.}
\begin{equation}
  \grad\cdot\vec{v}=\dpar{v_x}{x}+\dpar{v_y}{y}+\dpar{v_z}{z}
\end{equation}
é um escalar que mede o quanto o campo ``diverge'' de um ponto (fluxo por unidade
de volume).

\paragraph{Rotacional.}
\begin{equation}
  \grad\times\vec{v}=
  \begin{vmatrix}
    \uvec{x} & \uvec{y} & \uvec{z}\\[2pt]
    \dpar{}{x} & \dpar{}{y} & \dpar{}{z}\\[4pt]
    v_x & v_y & v_z
  \end{vmatrix}
\end{equation}
é um vetor que mede a circulação local (rotação) do campo.

\subsection{Regras de produto e segundas derivadas}

\add{As regras de produto úteis são:}
\begin{align}
  \grad(fg)&=f\,\grad g+g\,\grad f,\\
  \grad\cdot(f\vec{A})&=f\,(\grad\cdot\vec{A})+\vec{A}\cdot\grad f,\\
  \grad\times(f\vec{A})&=f\,(\grad\times\vec{A})+(\grad f)\times\vec{A},\\
  \grad\cdot(\vec{A}\times\vec{B})&=\vec{B}\cdot(\grad\times\vec{A})-\vec{A}\cdot(\grad\times\vec{B}).
\end{align}
As \textbf{segundas derivadas} relevantes:
\begin{align}
  \grad\cdot(\grad f)&=\nabla^2 f
     =\dpar{^2 f}{x^2}+\dpar{^2 f}{y^2}+\dpar{^2 f}{z^2}
     &&\text{(laplaciano)},\\
  \grad\times(\grad f)&=\vec{0} &&\text{(rotacional do gradiente)},\\
  \grad\cdot(\grad\times\vec{v})&=0 &&\text{(divergente do rotacional)},\\
  \grad\times(\grad\times\vec{v})&=\grad(\grad\cdot\vec{v})-\nabla^2\vec{v}.
\end{align}

\subsection{Operadores em coordenadas curvilíneas}

Este é o ponto em que o formalismo de índices rende fruto. Em um sistema de
coordenadas $\{x^i\}$ com base coordenada $\vec{e}_i=\partial\vec{r}/\partial x^i$
(em geral \emph{não} normalizada), define-se o \textbf{tensor métrico} e sua
inversa por
\begin{equation}
  g_{ij}=\langle\vec{e}_i,\vec{e}_j\rangle,\qquad
  g^{ik}g_{kj}=\delta^{i}{}_{j}.
\end{equation}
Como $\mathrm{d}f=(\partial_i f)\,\mathrm{d}x^i=\grad f\cdot\mathrm{d}\vec{l}$, e
elevando o índice com a métrica, o \textbf{gradiente} tem componentes
contravariantes
\begin{equation}
  (\grad f)^{\,j}=g^{ij}\,\partial_i f
  \qquad\Longrightarrow\qquad
  \grad f=\sum_{i,j} g^{ij}\,(\partial_i f)\,\vec{e}_j .
  \label{eq:gradcurv}
\end{equation}
A variação da base introduz os \textbf{símbolos de Christoffel}
\begin{equation}
  \dpar{\vec{e}_j}{x^k}=\Gamma^{\,i}{}_{jk}\,\vec{e}_i,
  \qquad
  \Gamma^{\,i}{}_{jk}=\left(\dpar{\vec{e}_j}{x^k}\right)^{i}.
\end{equation}
Com a \textbf{derivada covariante} $\nabla_i F^{j}=\partial_i F^{j}
+\Gamma^{\,j}{}_{ik}F^{k}$, o \textbf{divergente} e o \textbf{rotacional} ficam
\begin{align}
  \grad\cdot\vec{F}
    &=\nabla_i F^{i}=\partial_i F^{i}+\Gamma^{\,i}{}_{ik}F^{k},
    \label{eq:divcurv}\\
  (\grad\times\vec{F})^{i}
    &=\epsilon^{ijk}\,\nabla_j F_k
     =\epsilon^{ijk}\big(\partial_j F_k-\Gamma^{\,l}{}_{jk}F_l\big).
    \label{eq:rotcurv}
\end{align}
\begin{nota}
No rotacional, o termo com $\Gamma^{\,l}{}_{jk}$ \emph{se cancela}: os símbolos de
Christoffel são simétricos em $(j,k)$ e $\epsilon^{ijk}$ é antissimétrico, logo
$\epsilon^{ijk}\Gamma^{\,l}{}_{jk}=0$ e sobra $(\grad\times\vec{F})^{i}
=\epsilon^{ijk}\partial_j F_k$.
Além disso, vale a forma compacta do divergente (identidade de Voss--Weyl), com
$g=\det(g_{ij})$:
\begin{equation}
  \grad\cdot\vec{F}=\frac{1}{\sqrt{g}}\,\partial_i\!\big(\sqrt{g}\,F^{i}\big).
\end{equation}
\end{nota}

\paragraph{Caso ortogonal (fatores de escala).}
Quando a base é ortogonal, $g_{ij}=\operatorname{diag}(h_1^2,h_2^2,h_3^2)$ e os
versores normalizados são $\uvec{e}_i=\vec{e}_i/h_i$. Então
\begin{align}
  \grad f&=\sum_i \frac{1}{h_i}\dpar{f}{x^i}\,\uvec{e}_i,\\
  \grad\cdot\vec{F}&=\frac{1}{h_1h_2h_3}
     \left[\partial_1(h_2h_3F_1)+\partial_2(h_3h_1F_2)+\partial_3(h_1h_2F_3)\right],\\
  \nabla^2 f&=\frac{1}{h_1h_2h_3}
     \left[\partial_1\!\Big(\tfrac{h_2h_3}{h_1}\partial_1 f\Big)
          +\partial_2\!\Big(\tfrac{h_3h_1}{h_2}\partial_2 f\Big)
          +\partial_3\!\Big(\tfrac{h_1h_2}{h_3}\partial_3 f\Big)\right].
\end{align}

\paragraph{Exemplo (coordenadas polares/planas).}
Para $x=\rho\cos\varphi$, $y=\rho\sin\varphi$, obtém-se
$\vec{e}_\rho$ com $|\vec{e}_\rho|=1$ e $\vec{e}_\varphi$ com
$|\vec{e}_\varphi|=\rho$, isto é
\begin{equation}
  (g_{ij})=\begin{pmatrix}1&0\\0&\rho^2\end{pmatrix}
  \ \Longrightarrow\
  (g^{ij})=\begin{pmatrix}1&0\\0&\rho^{-2}\end{pmatrix},
  \qquad h_\rho=1,\ h_\varphi=\rho.
\end{equation}
Por \eqref{eq:gradcurv},
\(
  \grad f=\dpar{f}{\rho}\,\uvec{e}_\rho+\dfrac{1}{\rho}\dpar{f}{\varphi}\,\uvec{e}_\varphi .
\)
\add{Os únicos Christoffel não nulos são
$\Gamma^{\rho}{}_{\varphi\varphi}=-\rho$ e
$\Gamma^{\varphi}{}_{\rho\varphi}=\Gamma^{\varphi}{}_{\varphi\rho}=1/\rho$,}
consistentes com $\partial_\varphi\vec{e}_\varphi=-\rho\,\vec{e}_\rho$ etc.

\paragraph{Coordenadas esféricas.}
Para $\vec{r}=r\,\uvec{r}$, com $h_r=1$, $h_\theta=r$, $h_\varphi=r\sin\theta$,
\begin{equation}
  \grad=\uvec{r}\,\dpar{}{r}
       +\frac{\uvec{\theta}}{r}\,\dpar{}{\theta}
       +\frac{\uvec{\varphi}}{r\sin\theta}\,\dpar{}{\varphi},
\end{equation}
\add{e o laplaciano}
\begin{equation}
  \nabla^2 f=\frac{1}{r^2}\dpar{}{r}\!\Big(r^2\dpar{f}{r}\Big)
    +\frac{1}{r^2\sin\theta}\dpar{}{\theta}\!\Big(\sin\theta\,\dpar{f}{\theta}\Big)
    +\frac{1}{r^2\sin^2\theta}\dpar{^2 f}{\varphi^2}.
\end{equation}

\paragraph{Coordenadas cilíndricas.}
\add{Para $x=\rho\cos\varphi$, $y=\rho\sin\varphi$, $z=z$, tem-se
$h_\rho=1$, $h_\varphi=\rho$, $h_z=1$, donde
\(
  \nabla^2 f=\dfrac{1}{\rho}\partial_\rho(\rho\,\partial_\rho f)
  +\dfrac{1}{\rho^2}\partial_\varphi^2 f+\partial_z^2 f
\)
e o elemento de volume $\mathrm{d}\tau=\rho\,\mathrm{d}\rho\,\mathrm{d}\varphi\,\mathrm{d}z$.}

\subsection{Cálculo integral e teoremas fundamentais}

Definem-se a \textbf{integral de linha} $\int_{\mathcal{C}}\vec{v}\cdot\mathrm{d}\vec{l}$,
a \textbf{integral de superfície} (fluxo) $\int_{S}\vec{v}\cdot\mathrm{d}\vec{a}$
e a \textbf{integral de volume} $\int_{V}T\,\mathrm{d}\tau$. Valem os teoremas
fundamentais:
\begin{align}
  \int_{a}^{b}\!\Big(\dpar{f}{x}\Big)\mathrm{d}x&=f(b)-f(a)
     &&\text{(cálculo)},\\
  \int_{a}^{b}(\grad T)\cdot\mathrm{d}\vec{l}&=T(b)-T(a)
     &&\text{(gradiente)},\\
  \int_{V}(\grad\cdot\vec{F})\,\mathrm{d}\tau&=\oint_{\partial V}\vec{F}\cdot\mathrm{d}\vec{a}
     &&\text{(Gauss/divergência)},\\
  \int_{S}(\grad\times\vec{v})\cdot\mathrm{d}\vec{a}&=\oint_{\partial S}\vec{v}\cdot\mathrm{d}\vec{l}
     &&\text{(Stokes)}.
\end{align}
\add{Do teorema do gradiente seguem $\displaystyle\oint(\grad T)\cdot\mathrm{d}\vec{l}=0$
e a independência do caminho.}
\begin{nota}
Escrevendo o teorema de Gauss para $\vec{F}$, e as \emph{identidades de Green}
como consequência ($\phi\to$ escalar, $\chi\to$ escalar):
\begin{equation}
  \int_V\!\big[\phi\,\nabla^2\chi+\grad\phi\cdot\grad\chi\big]\mathrm{d}\tau
    =\oint_{S}\phi\,(\grad\chi)\cdot\mathrm{d}\vec{a}
    =\oint_S \phi\,\dpar{\chi}{n}\,\mathrm{d}a .
\end{equation}
\end{nota}

\subsection{A função delta de Dirac}

Em uma dimensão, $\delta(x)=0$ para $x\neq0$, com $\int\delta(x)\,\mathrm{d}x=1$;
rigorosamente, é uma \emph{distribuição} definida por
\begin{equation}
  \int f(x)\,\delta(x-a)\,\mathrm{d}x=f(a).
\end{equation}
Propriedades:
\begin{equation}
  \delta(ax)=\frac{1}{|a|}\,\delta(x),\qquad
  \delta\big(g(x)\big)=\sum_{n}\frac{\delta(x-x_n)}{|g'(x_n)|}
  \ \ \big(g(x_n)=0,\ g'(x_n)\neq0\big),
\end{equation}
\add{e a identidade de derivada $\displaystyle\int f(x)\,\delta'(x-x')\,\mathrm{d}x=-f'(x')$.}
Uma representação integral útil (Fourier):
\begin{equation}
  \delta(x-x')=\frac{1}{2\pi}\int_{-\infty}^{\infty}\mathrm{e}^{\,ik(x-x')}\,\mathrm{d}k
             =\int \mathrm{d}k\;\chi_k^{*}(x')\,\chi_k(x),
\end{equation}
\add{com $\chi_k(x)=\mathrm{e}^{ikx}/\sqrt{2\pi}$ (relação de completeza).}

Em três dimensões, $\delta^{3}(\rsep)=\delta(x-x')\,\delta(y-y')\,\delta(z-z')$,
com $\int f(\vec{r})\,\delta^{3}(\vec{r}-\vec{a})\,\mathrm{d}\tau=f(\vec{a})$.
O resultado central que resolve o ``paradoxo'' de $\grad\cdot(\usep/\Rsep^2)$ é
\begin{equation}
  \grad\cdot\!\left(\frac{\usep}{\Rsep^{2}}\right)=4\pi\,\delta^{3}(\rsep),
  \qquad
  \nabla^{2}\!\left(\frac{1}{\Rsep}\right)=-\,4\pi\,\delta^{3}(\rsep),
  \label{eq:delta3}
\end{equation}
\add{pois $\grad(1/\Rsep)=-\,\usep/\Rsep^{2}$.}

\subsection{Teoria dos campos vetoriais e potenciais}

\paragraph{Teorema de Helmholtz.} Um campo $\vec{C}$ que decai suficientemente
rápido no infinito decompõe-se em uma parte irrotacional e uma solenoidal,
\begin{equation}
  \vec{C}=\vec{C}_\perp+\vec{C}_\parallel,\qquad
  \grad\cdot\vec{C}_\perp=0,\qquad
  \grad\times\vec{C}_\parallel=\vec{0},
\end{equation}
e fica \emph{univocamente} determinado por seu divergente e seu rotacional. Uma
construção explícita é $\vec{C}(\vec{r})=\grad\times\vec{F}(\vec{r})-\grad\Omega(\vec{r})$
com
\begin{equation}
  \Omega(\vec{r})=\frac{1}{4\pi}\!\int \mathrm{d}^3r'\,
     \frac{\grad'\!\cdot\vec{C}(\vec{r}\,')}{|\vec{r}-\vec{r}\,'|},
  \qquad
  \vec{F}(\vec{r})=\frac{1}{4\pi}\!\int \mathrm{d}^3r'\,
     \frac{\grad'\!\times\vec{C}(\vec{r}\,')}{|\vec{r}-\vec{r}\,'|}.
\end{equation}
\begin{nota}
Esboço da prova: parte-se de
$\vec{C}(\vec{r})=\int \mathrm{d}^3r'\,\vec{C}(\vec{r}\,')\,\delta^3(\vec{r}-\vec{r}\,')$
e usa-se \eqref{eq:delta3} para trocar $\delta^3$ por
$-\tfrac{1}{4\pi}\nabla^2(1/|\vec{r}-\vec{r}\,'|)$. Aplicando
$\grad\times(\grad\times\vec{A})=\grad(\grad\cdot\vec{A})-\nabla^2\vec{A}$ e
integrando por partes (com $\grad(1/|\vec{r}-\vec{r}\,'|)=-\grad'(1/|\vec{r}-\vec{r}\,'|)$),
os termos de superfície vão a zero quando $\vec{C}$ decai no infinito.
\end{nota}

\paragraph{Potenciais.} Como consequência das segundas derivadas:
\begin{align}
  \grad\times\vec{F}=\vec{0}&\ \Rightarrow\ \vec{F}=-\grad V
  &&\text{(potencial escalar)},\\
  \grad\cdot\vec{F}=0&\ \Rightarrow\ \vec{F}=\grad\times\vec{A}
  &&\text{(potencial vetor)}.
\end{align}

%##############################################################################
\section{Eletrostática}
%##############################################################################

\subsection{Lei de Coulomb}

A força sobre uma carga $q$ devida a $N$ cargas $q_n$ em $\vec{r}_n$ é
\begin{equation}
  \vec{F}=\frac{q}{4\pi\eps}\sum_{n=1}^{N} q_n\,
          \frac{\vec{r}-\vec{r}_n}{|\vec{r}-\vec{r}_n|^{3}},
\end{equation}
\add{onde $\eps=8{,}85\times10^{-12}\ \mathrm{C^2/(N\,m^2)}$ é a permissividade do
vácuo.}

\subsection{Campo elétrico e distribuições contínuas}

Definindo $\vec{F}=q\,\vec{E}$, o \textbf{campo elétrico} de uma distribuição
contínua com densidade $\rho$ é
\begin{equation}
  \vec{E}(\vec{r})=\frac{1}{4\pi\eps}\!\int \mathrm{d}^3r'\,\rho(\vec{r}\,')\,
     \frac{\vec{r}-\vec{r}\,'}{|\vec{r}-\vec{r}\,'|^{3}}.
  \label{eq:Efield}
\end{equation}
\begin{nota}
Consistência com a força entre distribuições:
\begin{equation}
  \vec{F}=\frac{1}{4\pi\eps}\!\int\!\mathrm{d}^3r\!\int\!\mathrm{d}^3r'\,
    \rho(\vec{r})\,\rho(\vec{r}\,')\,
    \frac{\vec{r}-\vec{r}\,'}{|\vec{r}-\vec{r}\,'|^{3}}
  =\int \mathrm{d}^3r\,\rho(\vec{r})\,\vec{E}(\vec{r}).
\end{equation}
\end{nota}

\begin{obs}
\textbf{Comportamento sob escala.} Se $\rho_B(\lambda\vec{r})=\rho_A(\vec{r})$,
então de \eqref{eq:Efield} segue $\vec{E}_B(\lambda\vec{r})=\lambda\,\vec{E}(\vec{r})$
e o potencial escala como $\varphi_A(\vec{r})=\lambda^{-2}\varphi_B(\lambda\vec{r})$,
um teste dimensional útil.
\end{obs}

\subsection{Fluxo elétrico e lei de Gauss}

O carregamento total é $Q=\int_{S}\sigma\,\mathrm{d}S$ (superfície) ou
$Q=\int_{V}\rho\,\mathrm{d}V$ (volume). A \textbf{lei de Gauss} na forma integral,
\begin{equation}
  \Phi_E=\oint_{S}\vec{E}\cdot\mathrm{d}\vec{a}=\frac{Q_{\text{env}}}{\eps},
\end{equation}
combinada com o teorema da divergência,
$\oint_S\vec{E}\cdot\mathrm{d}\vec{a}=\int_V(\grad\cdot\vec{E})\,\mathrm{d}\tau$, e
com $Q_{\text{env}}=\int_V\rho\,\mathrm{d}\tau$, fornece a forma diferencial
\begin{equation}
  \boxed{\ \grad\cdot\vec{E}=\frac{\rho}{\eps}\ }
\end{equation}
— a primeira equação de Maxwell.
\add{A conservação da carga aparece já aqui: de $\vec{j}=\rho\,\vec{v}$ e
$\dot{Q}=\int \mathrm{d}^3x\,\partial_t\rho=\oint \mathrm{d}S\,\uvec{n}\cdot\vec{j}$
segue a equação de continuidade $\grad\cdot\vec{j}=-\,\partial_t\rho$
(no regime estático, $\partial_t\rho=0$).}

\subsection{Rotacional de \texorpdfstring{$\vec{E}$}{E} e potencial escalar}

A força de Coulomb é conservativa. Como
$\vec{E}=-\grad\varphi$ para o campo de uma carga pontual,
\begin{equation}
  q\!\int_{r_1}^{r_2}\!\mathrm{d}\vec{l}\cdot\vec{E}
  =-\,q\!\int\!\mathrm{d}\vec{l}\cdot\grad\varphi
  =-\,q\,[\varphi(r_2)-\varphi(r_1)],
\end{equation}
de modo que $W=q\oint\mathrm{d}\vec{l}\cdot\vec{E}=0$ e, por Stokes,
\begin{equation}
  \boxed{\ \grad\times\vec{E}=\vec{0}\ }.
\end{equation}
Por superposição, isto vale para qualquer distribuição. O \textbf{potencial
escalar eletrostático} é
\begin{equation}
  \varphi(\vec{r})=\frac{1}{4\pi\eps}\!\int \mathrm{d}^3r'\,
     \frac{\rho(\vec{r}\,')}{|\vec{r}-\vec{r}\,'|},
  \qquad \vec{E}=-\grad\varphi.
\end{equation}

\subsection{Equações de Poisson e Laplace}

Substituindo $\vec{E}=-\grad\varphi$ em $\grad\cdot\vec{E}=\rho/\eps$,
\begin{equation}
  \boxed{\ \nabla^{2}\varphi(\vec{r})=-\,\frac{\rho(\vec{r})}{\eps}\ }
  \qquad(\text{Poisson});\qquad
  \nabla^{2}\varphi=0\ \ (\text{Laplace, em regi\~oes sem carga}).
\end{equation}
\add{A solução $\varphi=\tfrac{1}{4\pi\eps}\int \rho/|\vec{r}-\vec{r}\,'|\,\mathrm{d}^3r'$
verifica Poisson por \eqref{eq:delta3}.}

\paragraph{Teorema de Earnshaw.} \add{Em uma região sem carga
($\nabla^2\varphi=0$), $\varphi$ não possui máximo nem mínimo locais no interior:
os extremos ocorrem necessariamente na fronteira. Consequência: não há
configuração de equilíbrio \emph{estável} puramente eletrostática para uma carga
de teste.}

\subsection{Condições de contorno}

Numa interface com densidade superficial $\sigma$, a componente normal de
$\vec{E}$ salta e a tangencial é contínua:
\begin{equation}
  \uvec{n}_2\cdot(\vec{E}_1-\vec{E}_2)=\frac{\sigma}{\eps},
\end{equation}
o que, em termos do potencial, dá a descontinuidade da derivada normal
\begin{equation}
  \left(\dpar{\varphi_2}{n_2}-\dpar{\varphi_1}{n_2}\right)\!\Bigg|_{S}
  =\textcolor{coradd}{-}\,\frac{\sigma}{\eps},
\end{equation}
\add{com $\varphi$ contínuo através de $S$. (O sinal segue de
$\vec{E}=-\grad\varphi$.)}

\begin{exemplo}[Potencial de uma linha carregada]
Considere um segmento sobre o eixo $z$, de $-L$ a $L$, com densidade linear
$\lambda$. Em um ponto $(x,y,z)$,
\begin{equation}
  \varphi(x,y,z)=\frac{1}{4\pi\eps}\!\int_{-L}^{L}
     \frac{\lambda\,\mathrm{d}z'}{\sqrt{x^2+y^2+(z-z')^2}}.
\end{equation}
Com $x^2+y^2=\rho^2$ e a troca $u=z'-z$ ($\mathrm{d}u=\mathrm{d}z'$),
\begin{equation}
  \varphi=\frac{\lambda}{4\pi\eps}\!\int_{-L-z}^{\,L-z}
     \frac{\mathrm{d}u}{\sqrt{\rho^2+u^2}}.
\end{equation}
Pondo $u=\rho\tan\theta$, $\mathrm{d}u=\rho\sec^2\theta\,\mathrm{d}\theta$ e
$\sqrt{\rho^2+u^2}=\rho\sec\theta$:
\begin{equation}
  \varphi=\frac{\lambda}{4\pi\eps}\!\int_{B}^{A}
     \frac{\rho\sec^2\theta}{\rho\sec\theta}\,\mathrm{d}\theta
  =\frac{\lambda}{4\pi\eps}\!\int_{B}^{A}\sec\theta\,\mathrm{d}\theta
  =\frac{\lambda}{4\pi\eps}\!\int_{B}^{A}\frac{\mathrm{d}\theta}{\cos\theta}.
\end{equation}
\begin{nota}
Como $\int\sec\theta\,\mathrm{d}\theta=\ln|\sec\theta+\tan\theta|$ e
$\sec\theta+\tan\theta=(\sqrt{\rho^2+u^2}+u)/\rho$, o resultado fechado é
\begin{equation}
  \varphi(\rho,z)=\frac{\lambda}{4\pi\eps}\,
    \ln\!\left[\frac{\sqrt{\rho^2+(L-z)^2}+(L-z)}
                    {\sqrt{\rho^2+(L+z)^2}-(L+z)}\right].
\end{equation}
No limite $L\to\infty$ recupera-se, a menos de constante,
$\varphi=-\tfrac{\lambda}{2\pi\eps}\ln\rho$, o fio infinito.
\end{nota}
\end{exemplo}

\begin{nota}
\textbf{Continuação pendente.} As páginas 19--22 do manuscrito não puderam ser
lidas nesta sessão (falha temporária na renderização das imagens). O documento
está estruturado para receber essas seções (provavelmente: energia eletrostática,
condutores e capacitância) sem alterar o que já foi digitado.
\end{nota}

%------------------------------------------------------------------ referências
\begin{thebibliography}{9}
\bibitem{zangwill}
  ZANGWILL, A. \textbf{Modern Electrodynamics}. Cambridge: Cambridge University
  Press, 2013.
\end{thebibliography}

\end{document}
